% ============================================================================
%   IX: A Tractable Descendant of AIXI for Autonomous Capital Compounding
%   Aldrich K. Wooden, Sr.  --  Zuup Innovation Lab / Visionblox LLC
%   Draft v0.1  --  May 2026
%
%   Build:
%     pdflatex IX_Whitepaper
%     bibtex   IX_Whitepaper
%     pdflatex IX_Whitepaper
%     pdflatex IX_Whitepaper
% ============================================================================

\documentclass[conference,10pt,letterpaper]{IEEEtran}
\IEEEoverridecommandlockouts

\usepackage[utf8]{inputenc}
\usepackage[T1]{fontenc}
\usepackage{cite}
\usepackage{amsmath,amssymb,amsfonts}
\usepackage{textcomp}
\usepackage{xcolor}
\usepackage{array}
\usepackage{tabularx}
\usepackage{booktabs}
\usepackage{url}
\usepackage[hidelinks]{hyperref}

% ---------------------------------------------------------------------------
%  Zuup epistemic markers (Verified / Plausible / Speculative)
% ---------------------------------------------------------------------------
\newcommand{\EV}{\textcolor{black}{\ensuremath{\checkmark}}\,}        % VERIFIED
\newcommand{\EP}{\textcolor{black}{\ensuremath{\circleddash}}\,}     % PLAUSIBLE
\newcommand{\ES}{\textcolor{black}{\ensuremath{\bigcirc}}\,}         % SPECULATIVE

% Declare the original Unicode markers as alternates (allows utf8 source too)
\DeclareUnicodeCharacter{2713}{\EV}
\DeclareUnicodeCharacter{25D0}{\EP}
\DeclareUnicodeCharacter{25EF}{\ES}
\DeclareUnicodeCharacter{00D7}{\ensuremath{\times}}
\DeclareUnicodeCharacter{2014}{---}
\DeclareUnicodeCharacter{2013}{--}
\DeclareUnicodeCharacter{2264}{\ensuremath{\leq}}
\DeclareUnicodeCharacter{2265}{\ensuremath{\geq}}
\DeclareUnicodeCharacter{0394}{\ensuremath{\Delta}}
\DeclareUnicodeCharacter{03B5}{\ensuremath{\varepsilon}}
\DeclareUnicodeCharacter{2248}{\ensuremath{\approx}}

% ===========================================================================
\begin{document}

\title{IX: A Tractable Descendant of AIXI \\ for Autonomous Capital Compounding}

\author{
  \IEEEauthorblockN{Aldrich K. Wooden, Sr., MBA}
  \IEEEauthorblockA{\textit{Zuup Innovation Lab} \quad \textit{Visionblox LLC} \\
                    Huntsville, AL, USA \\
                    Draft v0.1 --- May 2026}
}

\maketitle

% ---------------------------------------------------------------------------
\begin{abstract}
We introduce IX, a tractable architecture for an autonomous agent whose
objective is recursive self-financing (RSF): the systematic
deployment of recurring operator capital contributions and the
reinvestment of agent-generated returns toward a specified terminal
liquidity target. IX inherits its theoretical scaffolding
from Hutter's AIXI and its computational pattern from the MC-AIXI-CTW
approximation of Veness \textit{et al.}, substituting AIXI's uncomputable
universal prior with a domain-specific Bayesian world model over market
microstructure, and substituting its expectimax planner with policy-guided
Monte Carlo tree search under an explicit ruin-probability constraint. We
propose IX-Bench v0.1, an eight-family benchmark suite that grades agents on
returns, risk, recursion quality, and probability-of-target-achievement under
stated capital constraints. We outline a seven-phase implementation pipeline
mapped to USPTO prior-art clusters and open-source repository anchors. We
close with a non-optional treatment of empirical reality: a 2.5\%
monthly compound return (34.5\% annualized) sits in the top decile of
documented systematic returns---below Renaissance Medallion's
$\sim$40\% lifetime average but materially above commodity quant
benchmarks. The benchmark is structured to surface the gap between
assumed and achieved rather than obscure it.
\end{abstract}

\begin{IEEEkeywords}
AIXI, MC-AIXI-CTW, recursive self-financing, autonomous agents,
Monte Carlo tree search, Kelly criterion, algorithmic trading,
USPTO prior art
\end{IEEEkeywords}

% ===========================================================================
\section{Introduction}

A recursively self-financing agent is a closed-loop system that takes
positions in markets, reinvests realized gains, updates its own policy from
realized outcomes, and does so without human intervention beyond strategic
configuration. The construct composes three problems usually studied in
isolation: sequential decision under uncertainty, online learning under
non-stationarity, and capital allocation under irreversibility. The first
has a clean theoretical ceiling (AIXI). The second is the subject of
decades of online-learning literature. The third is the source of practical
ruin for most agents that succeed in the first two.

The motivating scenario for IX is a recurring operator capital
contribution of \$5{,}000 per month, deployed by the agent under a 2.5\%
monthly target return (34.5\% annualized), toward a terminal liquidity
target of \$10M after taxes. Under the stated parameters, the
deterministic time-to-target is approximately 159 months ($\approx$13.3
years) pre-tax and 173--184 months ($\approx$14.4--15.3 years) after
taxes depending on the realized short-term/long-term gains mix---a
tractable horizon, and the 2.5\%/month return assumption sits within
the historically documented distribution of elite systematic
performance, though sustained achievement across the full horizon is
not commodity. Epistemic markers are used throughout: \EV{}VERIFIED,
\EP{}PLAUSIBLE, \ES{}SPECULATIVE.

% ===========================================================================
\section{Theoretical Foundations}

\subsection{The AIXI Ceiling}

AIXI defines the optimal action for an agent in any computable environment
as the action that maximizes expected future reward, where the expectation
is taken over all computable environment hypotheses weighted by their
Kolmogorov complexity~\cite{hutter2005universal}. Concretely, the optimal
action at time $t$ is the argmax over $a_t$ of the expectimax expansion
across future observations and rewards, with the inner mixture over programs
$q$ running on a universal Turing machine $U$, each weighted by
$2^{-\ell(q)}$ where $\ell(q)$ is program length~\cite{solomonoff1964formal}.
\EV{}AIXI is asymptotically optimal across the class of all computable
environments. \EV{}AIXI is uncomputable: Kolmogorov complexity is not a
computable function.

\subsection{The MC-AIXI-CTW Tractability Pattern}

Veness, Ng, Hutter, Uther, and Silver~\cite{veness2011monte} demonstrated
that AIXI's two-slot structure survives two surgical substitutions. The
world-model slot becomes Context Tree Weighting (CTW): a Bayesian mixture
over variable-order Markov models bounded by depth $D$, which preserves the
mixture-of-experts character of Solomonoff induction inside a tractable model
class. The planner slot becomes Monte Carlo tree search (MCTS): expectimax
is replaced by stochastic rollouts that concentrate compute on promising
branches~\cite{silver2016mastering}. The composed system is computable on
commodity hardware and learned non-trivial domains (Pac-Man, Kuhn poker,
TicTacToe) from scratch.

\subsection{The IX Substitution}

IX preserves the two-slot pattern but specializes each slot for the RSF
domain. The world model is a hybrid: a regime-aware state-space model over
price and volume, an on-chain liquidity and orderflow model, and a finite
Bayesian mixture over candidate trading hypotheses whose weights are updated
by realized PnL---preserving the AIXI-style discipline of crediting theories
that predict correctly and shrinking those that do not. The planner is MCTS
guided by a policy network trained via PPO~\cite{schulman2017proximal} or
its successor, with rollouts that evaluate not only expected return but the
realized empirical ruin probability under each candidate action. The
objective is therefore constrained:
%
\begin{equation}
\pi^{*} = \arg\max_{\pi \in \Pi}\; \mathbb{E}_{\pi}[R]
\quad \text{subject to} \quad
\Pr(\text{ruin} \mid \pi) \leq \varepsilon.
\label{eq:objective}
\end{equation}
%
\EP{}The constrained variant is essential when capital is irreversibly lost
on failure; the unconstrained AIXI objective is silently incompatible with
leveraged or concentrated betting.

% ===========================================================================
\section{Architecture: The Seven Layers}

IX decomposes into seven layers, each addressing a single concern.
Inter-layer interfaces are explicit data contracts so that any single layer
can be reimplemented or replaced without disturbing the rest
(Table~\ref{tab:architecture}).

\begin{table}[!t]
\caption{IX Seven-Layer Architecture}
\label{tab:architecture}
\centering
\renewcommand{\arraystretch}{1.15}
\begin{tabularx}{\columnwidth}{@{}c l X@{}}
\toprule
\textbf{Phase} & \textbf{Layer} & \textbf{Concern} \\
\midrule
0 & Substrate                  & Data ingestion; on/off-chain unified feed; latency-aware reconciliation \\
1 & World Model                & Bayesian posterior over market regimes and competing hypotheses \\
2 & Strategy                   & Policy generation under complexity-penalty regularization \\
3 & Execution                  & Routing with information-leakage cost internalized \\
4 & Risk                       & Recursive Kelly damped by self-calibration error; ruin-bounded sizing \\
5 & Compounding                & Real-time tax-aware reinvestment scheduling \\
6 & Self-Improvement           & Meta-learning grounded in live out-of-sample PnL \\
7 & Compliance \& Tax          & After-tax reward signal; 1099-DA/8949 audit-trail by construction \\
\bottomrule
\end{tabularx}
\end{table}

Two architectural choices distinguish IX from a generic deep-RL trader.
First, the risk layer implements recursive Kelly~\cite{kelly1956new}:
position size is the Kelly fraction multiplied by a damping coefficient that
is a function of the agent's own recent calibration error, measured as the
empirical Brier score on its forward predictions---an overconfident agent
shrinks its own bet size automatically. Second, the self-improvement layer
treats meta-loss as live realized PnL on small bounded capital, not
backtested PnL, grounding self-modification in irreversible feedback rather
than self-deception about a model that overfit the test set.

% ===========================================================================
\section{Benchmark Framework: IX-Bench v0.1}

IX-Bench evaluates candidate agents across eight metric families. Each
family resolves to a 0--100 sub-score; sub-scores compose into a single IX
Capability Score (IX-CS) with weights configurable per operational phase.
Family~8---Probabilistic Target Achievement---is the family most often
omitted by performance frameworks and is non-optional here
(Table~\ref{tab:benchmark}).

\begin{table}[!t]
\caption{IX-Bench v0.1 Metric Families}
\label{tab:benchmark}
\centering
\renewcommand{\arraystretch}{1.15}
\begin{tabularx}{\columnwidth}{@{}c l X@{}}
\toprule
\textbf{\#} & \textbf{Family} & \textbf{Representative Metrics} \\
\midrule
1 & Return              & CAGR, TWR, IRR \\
2 & Risk                & MaxDD, DD Duration, VaR/CVaR, Risk-of-Ruin, Sharpe, Sortino, Calmar \\
3 & Strategy Quality    & Win Rate, Avg Win/Loss, Profit Factor, Recovery Factor, Tail Ratio, Info Ratio vs BTC-HODL \\
4 & Agent Performance   & Decision Latency, Brier Calibration, Tool Success, Hallucination Rate, OPE, Strategy Diversity \\
5 & Capital Efficiency  & Utilization, Reinvestment Velocity, Slippage, Fee Ratio \\
6 & Recursion (RSF)     & Gen-over-Gen Sharpe $\Delta$, OOS Degradation, Walk-Forward, Regime Detection, Alpha Half-Life \\
7 & Compliance/Ops      & After-Tax Net, 1099-DA/8949 Accuracy~\cite{irs2024digital}, AML FP Rate, Audit Coverage, Custody Risk \\
8 & Probabilistic Target & $\Pr(\$10\text{M after-tax by m175})$, $\Pr(\text{cumulative-return-negative at horizon})$, $\mathbb{E}[\text{time to target}]$, return-distribution quantiles, monthly-return drawdown distribution, model misspec sensitivity \\
\bottomrule
\end{tabularx}
\end{table}

% ===========================================================================
\section{Implementation Pipeline}

Each phase pairs a USPTO prior-art cluster---a region of patent space
densely covered by existing claims---with a non-obvious gap and an
open-source repository anchor that minimizes the rebuild surface
(Table~\ref{tab:pipeline}). Build stack: Claude Code for primary
implementation; Cursor for IDE iteration on Phases 1--2; Gemini for
second-pass verification of Kelly derivations, ruin-probability arithmetic,
and patent-claim novelty against the USPTO record; Claude (interactive
thread) for architecture and document iteration.

\begin{table*}[!t]
\caption{IX Implementation Pipeline: USPTO Prior-Art Clusters and Open-Source Anchors}
\label{tab:pipeline}
\centering
\renewcommand{\arraystretch}{1.2}
\begin{tabularx}{\textwidth}{@{}l l X X@{}}
\toprule
\textbf{Phase} & \textbf{USPTO Cluster (CPC)} & \textbf{Non-Obvious Gap} & \textbf{Open-Source Anchors} \\
\midrule
0~--~Substrate         & G06Q40/04             & Cross-feed disagreement treated as alpha signal, not noise to average away                       & ccxt, web3.py / ethers.js, clickhouse / arctic, polars / duckdb \\
1~--~World Model       & G06N5/04, G06N3/08    & Explicit Bayesian posterior over competing market hypotheses; disagreement as exploration signal & hmmlearn, pytorch-forecasting, mlfinlab, vectorbt, tslearn \\
2~--~Strategy          & G06N3/126, G06Q40/00  & Strategy synthesis under explicit Kolmogorov-complexity penalty as generalization regularizer    & DEAP, PyGAD, stable-baselines3, Ray RLlib, FinRL \\
3~--~Execution         & G06Q40/04             & Routing that internalizes the agent's own information-leakage cost; slippage treated as endogenous & hummingbot, jesse, freqtrade, 1inch SDK, cowswap \\
4~--~Risk              & G06Q40/06             & Recursive Kelly damped by self-calibration error (Brier-score-driven sizing damping)             & empyrical, pyfolio, riskfolio-lib, cvxpy, quantstats \\
5~--~Compounding       & G06Q40/06             & Real-time pricing of marginal compounding opportunity vs.\ marginal tax cost                     & defi-sdk, aave-v3-sdk, backtesting.py \\
6~--~Self-Improvement  & G06N3/08, G06N20/00   & Meta-loss grounded in live OOS PnL on bounded real capital, not backtest                         & Optuna, Ray Tune, NNI, MAML reference implementations \\
7~--~Compliance        & G06Q40/12             & RL reward signal = after-tax net realized, not gross PnL (1099-DA-aware)                         & rotki, open IRS-form generators \\
\bottomrule
\end{tabularx}
\end{table*}

% ===========================================================================
\section{Risk, Reality, and Honest Limitations}

Any whitepaper that omits the following is engaged in rhetorical performance
rather than research.

\subsection{The 2.5\%-Monthly Calibration}

\EV{}A 2.5\% monthly compound return corresponds to a 34.5\%
annualized return. This sits below Renaissance Medallion's lifetime
average of $\sim$40\% annually net of fees but materially above
commodity quant benchmarks: top-quartile systematic CTAs have
historically realized 15--25\% annualized; top-decile hedge funds
20--30\%; the S\&P~500 long-run $\sim$10\%. \EP{}The 2.5\%/month
assumption therefore sits in the top decile of documented systematic
performance over multi-year horizons---aggressive but not at the
boundary of the historical record. \EP{}Achievability is more
plausible at the \$5{,}000/month scale than at fund scale: small-capital
strategies can exploit microstructure regimes that do not scale, can
deploy leverage flexibly without meaningful market impact, and face less
alpha decay from public-record competition. \ES{}Nonetheless, sustained
2.5\%/month average across the full 159-month pre-tax horizon---through
multiple regime transitions, asset rotations, and bear markets---is
not commodity. Even Medallion has had down years. The architectural
defenses against return-rate erosion are: regime-adaptive strategy
switching (Phase~1's mixture posterior); recursive Kelly with
self-calibration damping (Phase~4) to compress drawdowns when
overconfidence is detected; and tax-aware holding-period optimization
(Phase~7) whose value grows in absolute terms as the return rate falls.
The agent's responsibility under IX-Bench is to report the realized
rolling return distribution honestly, not to claim deterministic
achievement of the geometric mean.

\subsection{Regulatory and Tax Surface}

\EV{}Autonomous trading by U.S.\ persons is subject to CFTC, SEC, and
FinCEN regimes depending on instrument class. \EV{}Digital asset gains
are reported on Form 8949 and Schedule D, with broker reporting under
Form 1099-DA effective from 2026~\cite{irs2024digital}. \EV{}Tax drag
is material to the \$10M-after-tax target: assuming the top U.S.\
federal marginal bracket plus Alabama state income tax (5\%), gains
treated as long-term ($\geq$1 year holding) face an effective rate of
$\sim$29\% (20\% federal LTCG + 3.8\% NIIT + 5\% state), while gains
treated as short-term face $\sim$46\% (37\% federal ordinary + 3.8\%
NIIT + 5\% state). The corresponding gross targets are \$14.1M and
\$18.5M respectively, translating to 173 months at the long-term
extreme and 184 months at the short-term extreme, against the
159-month pre-tax horizon. Active monthly rebalancing pushes most
realized gains into the short-term regime by default; the Phase~7
compliance layer must therefore actively manage holding-period
optimization to materially compress the time-to-target. Note that the
tax-drag penalty in absolute months grows as the assumed return rate
falls: an 11-month tax penalty at 2.5\%/month corresponds to the same
dollar tax burden as the 5--12-month penalty at 5\%/month, but is more
costly in calendar time because the underlying compounding is slower. \EP{}Operating an autonomous
agent on behalf of any party other than the operator introduces
fiduciary and registration questions whose resolution depends on
jurisdiction and on whether the operator is acting as an investment
adviser. The compliance layer (Phase~7) is non-negotiable and is
therefore architecturally co-equal with the return-generating layers.

% ===========================================================================
\section{Conclusion and Future Work}

IX is a research program, not a product. Its contribution is threefold: a
clean theoretical lineage from AIXI through MC-AIXI-CTW to a domain-specific
tractable agent; a benchmark framework that grades agents probabilistically
and refuses to obscure the gap between target and feasible; and a
seven-phase pipeline mapped to USPTO prior-art clusters and open-source
anchors. Future work: derivation of the constrained-MCTS variant with
formal ruin-probability bounds; empirical calibration of CTW depth for
crypto microstructure; formal definition of the IX-RSI metric (Recursive
Sharpe Improvement); and small-capital live trials with hard ruin-probability
ceilings to populate Family~8 of IX-Bench with empirical posteriors.

% ===========================================================================
\bibliographystyle{IEEEtran}
\bibliography{Ref}

\end{document}
